\chapter{The Package Manager}\label{ch:package-manager}

\index{package manager}\index{clat!package manager}

In the previous chapter, we learned how to split our own code across modules and import them.
But programming does not happen in isolation.
Eventually, you will want to use a library someone else wrote---a JSON schema validator, a CLI argument parser, a color formatting toolkit.
That is where the package manager comes in.

Lattice's package manager is built directly into \texttt{clat}\index{clat}, the same binary you use to run programs.
There is no separate tool to install, no Node.js-style global package manager that updates independently from the language.
Everything lives in one place: your manifest, your lock file, and a \filename{lat\_modules/}\index{lat\_modules} directory that works much like you would expect if you have used \texttt{npm} or \texttt{cargo}.

Let's start by creating a project.

%% ───────────────────────────────────────────────────────────
\section{Initializing a Project with \texttt{clat init}}\label{sec:clat-init}
\index{clat init}\index{lattice.toml!creating}

Every managed Lattice project begins with a \filename{lattice.toml}\index{lattice.toml} file.
You can create one by hand, but the fastest way is:

\begin{lstlisting}[language=bash, numbers=none]
mkdir weather-app && cd weather-app
clat init
\end{lstlisting}

Output:

\begin{lstlisting}[language=bash, numbers=none]
Created lattice.toml
\end{lstlisting}

Open the generated file and you will see:

\begin{lstlisting}[language=bash, numbers=none, caption={A freshly generated lattice.toml}]
[package]
name = "weather-app"
version = "0.1.0"
\end{lstlisting}

Three things to notice:

\begin{enumerate}
  \item The \texttt{name} is derived from the current directory name.
  \item The \texttt{version} starts at \texttt{0.1.0}, following semantic versioning conventions.
  \item There is no explicit \texttt{entry} field---when omitted, Lattice defaults to \filename{main.lat} as the entry point.
\end{enumerate}

If you run \lstinline|clat init| in a directory that already has a \filename{lattice.toml}, it refuses and prints an error.
It will never overwrite an existing manifest.

\begin{tipbox}{Custom Entry Points}
If your project's main file is not called \filename{main.lat}, add an \texttt{entry} field to the \texttt{[package]} section:

\begin{lstlisting}[language=bash, numbers=none]
[package]
name = "weather-app"
version = "0.1.0"
entry = "src/app.lat"
\end{lstlisting}

This tells both the package manager and importers where to find the project's primary module.
\end{tipbox}

%% ───────────────────────────────────────────────────────────
\section{The \texttt{lattice.toml} Manifest}\label{sec:lattice-toml}
\index{lattice.toml}\index{manifest}

The \filename{lattice.toml} file is the single source of truth for your project's metadata and dependencies.
It uses TOML\index{TOML} syntax---the same format Rust uses for \texttt{Cargo.toml} and Python uses for \texttt{pyproject.toml}.

Here is a complete manifest for a project with dependencies:

\begin{lstlisting}[language=bash, numbers=none, caption={A lattice.toml with metadata and dependencies}]
[package]
name = "weather-app"
version = "1.2.0"
description = "A weather monitoring application"
license = "MIT"
entry = "main.lat"

[dependencies]
http-client = "^1.0.0"
json-schema = "~0.3.2"
cli-args = "2.1.0"
color-fmt = "*"
\end{lstlisting}

\subsection{The \texttt{[package]} Section}\label{subsec:package-section}
\index{lattice.toml!package section}

The \texttt{[package]} table contains your project's metadata:

\begin{table}[h]
\centering
\caption{Fields in the \texttt{[package]} section}\label{tab:package-fields}
\begin{tabular}{@{}lll@{}}
\toprule
\textbf{Field} & \textbf{Required} & \textbf{Description} \\
\midrule
\texttt{name}        & Yes & The package name (derived from directory by \texttt{clat init}) \\
\texttt{version}     & Yes & Semantic version string (e.g., \texttt{"1.2.3"}) \\
\texttt{description} & No  & A short description of the package \\
\texttt{license}     & No  & SPDX license identifier (e.g., \texttt{"MIT"}, \texttt{"Apache-2.0"}) \\
\texttt{entry}       & No  & Entry point file (defaults to \texttt{"main.lat"}) \\
\bottomrule
\end{tabular}
\end{table}

\subsection{The \texttt{[dependencies]} Section}\label{subsec:dependencies-section}
\index{lattice.toml!dependencies}\index{dependencies}

The \texttt{[dependencies]} table maps package names to version constraints.
Each key is a package name, and each value is a string specifying which versions are acceptable:

\begin{lstlisting}[language=bash, numbers=none]
[dependencies]
http-client = "^1.0.0"
json-schema = "~0.3.2"
cli-args = "2.1.0"
color-fmt = "*"
\end{lstlisting}

We will break down what those version strings mean in \cref{sec:semver}.

%% ───────────────────────────────────────────────────────────
\section{Adding, Removing, and Installing Packages}\label{sec:add-remove-install}
\index{clat add}\index{clat remove}\index{clat install}

The package manager provides three commands for managing dependencies.

\subsection{\texttt{clat add}}\label{subsec:clat-add}
\index{clat add}

To add a dependency to your project:

\begin{lstlisting}[language=bash, numbers=none]
clat add http-client ^1.0.0
\end{lstlisting}

Output:

\begin{lstlisting}[language=bash, numbers=none]
Added http-client@^1.0.0 to dependencies
  Installing http-client@^1.0.0... ok

All dependencies installed.
\end{lstlisting}

This does two things:

\begin{enumerate}
  \item Adds \lstinline|http-client = "^1.0.0"| to the \texttt{[dependencies]} section of your \filename{lattice.toml}.
  \item Runs \texttt{clat install} to fetch and install the package.
\end{enumerate}

If you omit the version constraint, it defaults to \lstinline|"*"| (any version):

\begin{lstlisting}[language=bash, numbers=none]
clat add color-fmt
\end{lstlisting}

If the package is already a dependency, \texttt{clat add} updates its version constraint instead of adding a duplicate.

\begin{notebox}{No \texttt{lattice.toml} Yet?}
If you run \texttt{clat add} before running \texttt{clat init}, the package manager creates a minimal \filename{lattice.toml} for you automatically, deriving the package name from the current directory.
You do not need a separate initialization step if you start by adding dependencies.
\end{notebox}

\subsection{\texttt{clat install}}\label{subsec:clat-install}
\index{clat install}

To install all dependencies listed in your manifest:

\begin{lstlisting}[language=bash, numbers=none]
clat install
\end{lstlisting}

Output:

\begin{lstlisting}[language=bash, numbers=none]
  Installing http-client@^1.0.0... ok
  Installing json-schema@~0.3.2... ok
  Installing cli-args@2.1.0... ok
  Installing color-fmt@*... ok

All dependencies installed.
\end{lstlisting}

For each dependency, the installer follows this resolution strategy:

\begin{enumerate}
  \item \textbf{Check \texttt{lat\_modules/}.}\index{lat\_modules}
  If the package directory already exists, verify that the installed version satisfies the constraint.
  If it does, skip the download.

  \item \textbf{Check the lock file.}
  If \filename{lattice.lock}\index{lattice.lock} exists, prefer the locked version (as long as it still satisfies the manifest constraint).
  This ensures reproducible builds.

  \item \textbf{Check the local file registry.}
  If the \texttt{LATTICE\_REGISTRY}\index{LATTICE\_REGISTRY} environment variable is set to a \lstinline|file://| URL, try copying the package from that directory.

  \item \textbf{Fetch from the HTTP registry.}\index{registry}
  Query the registry for available versions, find the best match, download it to the global cache, and copy it into \filename{lat\_modules/}.
\end{enumerate}

After all dependencies are installed, \texttt{clat install} builds a dependency graph\index{dependency graph} and checks for circular dependencies\index{circular dependencies}.
If a cycle is detected, it reports the full cycle path and fails:

\begin{lstlisting}[language=bash, numbers=none]
error: circular dependency detected: alpha -> beta -> gamma -> alpha
\end{lstlisting}

Finally, it writes (or updates) the \filename{lattice.lock} file.

\subsection{\texttt{clat remove}}\label{subsec:clat-remove}
\index{clat remove}

To remove a dependency:

\begin{lstlisting}[language=bash, numbers=none]
clat remove color-fmt
\end{lstlisting}

Output:

\begin{lstlisting}[language=bash, numbers=none]
Removed color-fmt
\end{lstlisting}

This does three things:

\begin{enumerate}
  \item Removes the package from \filename{lattice.toml}.
  \item Deletes \filename{lat\_modules/color-fmt/} from disk.
  \item Removes the entry from \filename{lattice.lock}.
\end{enumerate}

If the package is not listed as a dependency, \texttt{clat remove} prints an error and does nothing.

%% ───────────────────────────────────────────────────────────
\section{Semantic Versioning and Dependency Resolution}\label{sec:semver}
\index{semantic versioning}\index{semver}

Lattice uses semantic versioning (semver) for all version strings.
A version has three numeric components:

\begin{defnbox}{Semantic Version}
\index{semver!definition}
A version string \texttt{MAJOR.MINOR.PATCH} communicates intent:

\begin{itemize}[nosep]
  \item \textbf{MAJOR} --- incremented for breaking, incompatible changes.
  \item \textbf{MINOR} --- incremented for backward-compatible new features.
  \item \textbf{PATCH} --- incremented for backward-compatible bug fixes.
\end{itemize}

For example, \texttt{2.3.1} means major version 2, minor version 3, patch level 1.
\end{defnbox}

\subsection{Version Constraints}\label{subsec:version-constraints}
\index{semver!constraints}\index{version constraints}

When you list a dependency in \filename{lattice.toml}, the version string is a \emph{constraint}, not an exact version.
Lattice supports six constraint types:

\begin{table}[h]
\centering
\caption{Semver constraint types}\label{tab:semver-constraints}
\begin{tabular}{@{}llp{7cm}@{}}
\toprule
\textbf{Syntax} & \textbf{Name} & \textbf{Meaning} \\
\midrule
\texttt{"*"}       & Wildcard & Any version. Useful during early development. \\
\texttt{"1.2.3"}   & Exact    & Exactly version 1.2.3 and nothing else. \\
\texttt{"\^{}1.2.3"} & Caret (Compatible) & Same major version, at least 1.2.3. Matches \texttt{1.2.3}, \texttt{1.3.0}, \texttt{1.99.0}, but not \texttt{2.0.0}. \\
\texttt{"\~{}1.2.3"} & Tilde (Approximate) & Same major \emph{and} minor version, at least patch 3. Matches \texttt{1.2.3}, \texttt{1.2.9}, but not \texttt{1.3.0}. \\
\texttt{">=1.2.3"} & Minimum  & Version 1.2.3 or newer. No upper bound. \\
\texttt{"<=1.2.3"} & Maximum  & Version 1.2.3 or older. No lower bound. \\
\bottomrule
\end{tabular}
\end{table}

\subsection{Caret vs.\ Tilde}\label{subsec:caret-vs-tilde}
\index{semver!caret}\index{semver!tilde}

The most commonly used constraints are caret (\texttt{\^{}}) and tilde (\texttt{\~{}}).
Here is how they differ in practice:

\begin{lstlisting}[language=bash, numbers=none, caption={Caret vs.\ tilde constraints}]
# Caret: ^1.2.3
#   Allows: 1.2.3, 1.2.4, 1.3.0, 1.99.0
#   Blocks: 2.0.0 (different major)

# Tilde: ~1.2.3
#   Allows: 1.2.3, 1.2.4, 1.2.99
#   Blocks: 1.3.0 (different minor)
\end{lstlisting}

\textbf{When to use caret:} You trust the package author to maintain backward compatibility within a major version.
This is the default choice for most dependencies.

\textbf{When to use tilde:} You want only bug fixes for a specific feature set.
This is more conservative and useful when a minor version introduced a behavior you depend on.

\begin{tipbox}{Default to Caret}
If you are unsure which constraint to use, start with caret (\texttt{\^{}}).
It gives package authors room to ship new features without breaking your code, while still protecting you from major-version breaking changes.
\end{tipbox}

\subsection{How Resolution Works}\label{subsec:resolution-algorithm}
\index{dependency resolution}

When \texttt{clat install} needs to resolve a version constraint, it follows this process:

\begin{enumerate}
  \item \textbf{Query the registry}\index{registry} for all published versions of the package.
  The registry endpoint is \texttt{GET /packages/\textlangle{}name\textrangle{}/versions}, which returns a JSON array.

  \item \textbf{Filter} the version list against the constraint using \lstinline|pkg_semver_satisfies()| (implemented in \filename{src/package.c}).

  \item \textbf{Select the highest} matching version.
  Lattice always prefers the newest compatible version, following the common ``maximal version'' strategy.
\end{enumerate}

For example, if the registry has versions \texttt{[1.0.0, 1.1.0, 1.2.0, 2.0.0]} and your constraint is \texttt{\^{}1.0.0}, the resolver picks \texttt{1.2.0}---the highest version that shares the same major version.

%% ───────────────────────────────────────────────────────────
\section{The Lock File and Reproducible Builds}\label{sec:lock-file}
\index{lattice.lock}\index{lock file}\index{reproducible builds}

Every time \texttt{clat install} succeeds, it writes a \filename{lattice.lock} file.
This file records the \emph{exact} versions that were installed, so that future installs produce identical results.

\subsection{What the Lock File Looks Like}\label{subsec:lock-format}
\index{lattice.lock!format}

\begin{lstlisting}[language=bash, numbers=none, caption={A typical lattice.lock}]
# This file is auto-generated by clat. Do not edit manually.

[[package]]
name = "http-client"
version = "1.2.0"
source = "registry"
checksum = ""

[[package]]
name = "json-schema"
version = "0.3.5"
source = "registry"
checksum = ""

[[package]]
name = "cli-args"
version = "2.1.0"
source = "local"
checksum = ""
\end{lstlisting}

Each \texttt{[[package]]}\index{lattice.lock!package entry} entry records:

\begin{itemize}[nosep]
  \item \textbf{\texttt{name}} --- the package name.
  \item \textbf{\texttt{version}} --- the exact resolved version.
  \item \textbf{\texttt{source}} --- where the package came from: \texttt{"registry"} (HTTP registry), \texttt{"local"} (already in \filename{lat\_modules/}), or \texttt{"path"} (file-based registry).
  \item \textbf{\texttt{checksum}} --- a SHA-256 hash for integrity verification (reserved for future use).
\end{itemize}

\subsection{How the Lock File Ensures Reproducibility}\label{subsec:lock-reproducibility}

When \texttt{clat install} finds an existing \filename{lattice.lock}, it prefers the locked version for each dependency---\emph{as long as it still satisfies the constraint in \filename{lattice.toml}}.
This means:

\begin{itemize}
  \item \textbf{Same machine, next week:} Running \lstinline|clat install| again installs the same versions, even if newer versions have been published.

  \item \textbf{Different machine, same lock file:} A teammate clones your repository, runs \lstinline|clat install|, and gets the exact same dependency versions.

  \item \textbf{Updating the constraint:} If you change a version constraint in \filename{lattice.toml} so that the locked version no longer satisfies it, the next \lstinline|clat install| fetches a new version and updates the lock file.
\end{itemize}

\begin{warnbox}{Commit Your Lock File}
Always commit \filename{lattice.lock} to version control.
This is how you guarantee reproducible builds across machines and over time.
The \filename{lat\_modules/} directory, on the other hand, should be in your \filename{.gitignore}---it can always be reconstructed from the lock file.
\end{warnbox}

\subsection{When Lock Entries Are Updated}\label{subsec:lock-updates}

The lock file is regenerated in these situations:

\begin{enumerate}
  \item \texttt{clat install} runs and resolves new or changed dependencies.
  \item \texttt{clat add} adds a new dependency and triggers an install.
  \item \texttt{clat remove} removes a dependency and strips its entry from the lock file.
\end{enumerate}

%% ───────────────────────────────────────────────────────────
\section{The Package Registry and Cache}\label{sec:registry-cache}
\index{registry}\index{package cache}

\subsection{The Lattice Registry}\label{subsec:registry}
\index{registry!HTTP}

Lattice's default package registry\index{registry!default} lives at:

\begin{lstlisting}[language=bash, numbers=none]
https://registry.lattice-lang.org/v1
\end{lstlisting}

The package manager communicates with the registry using two HTTP endpoints:

\begin{table}[h]
\centering
\caption{Registry API endpoints}\label{tab:registry-endpoints}
\begin{tabular}{@{}ll@{}}
\toprule
\textbf{Endpoint} & \textbf{Purpose} \\
\midrule
\texttt{GET /packages/\textlangle{}name\textrangle{}/versions} & List available versions (JSON) \\
\texttt{GET /packages/\textlangle{}name\textrangle{}/\textlangle{}version\textrangle{}} & Download package source \\
\bottomrule
\end{tabular}
\end{table}

The versions endpoint returns JSON like \lstinline|\{"versions":["1.0.0","1.1.0","2.0.0"]\}|.
The download endpoint returns either a single \texttt{main.lat} file or a TOML bundle, depending on the package.

\subsection{Overriding the Registry}\label{subsec:registry-override}
\index{LATTICE\_REGISTRY}\index{registry!override}

You can point the package manager at a different registry using the \texttt{LATTICE\_REGISTRY} environment variable:

\begin{lstlisting}[language=bash, numbers=none]
# Use a private registry
export LATTICE_REGISTRY=https://packages.mycompany.com/v1
clat install

# Use a local directory as a registry (great for testing)
export LATTICE_REGISTRY=file:///home/user/my-packages
clat install
\end{lstlisting}

The \lstinline|file://| protocol is particularly useful for development.
If you set \lstinline|LATTICE_REGISTRY=file:///path/to/packages|, the package manager copies packages directly from that directory into \filename{lat\_modules/} without any HTTP requests.

\subsection{The Global Package Cache}\label{subsec:global-cache}
\index{package cache!global}\index{\textasciitilde/.lattice/packages}

When a package is downloaded from the registry, it is stored in a global cache before being copied into your project's \filename{lat\_modules/}:

\begin{lstlisting}[language=bash, numbers=none]
~/.lattice/packages/<name>/<version>/
\end{lstlisting}

For example, \texttt{http-client} version \texttt{1.2.0} would be cached at:

\begin{lstlisting}[language=bash, numbers=none]
~/.lattice/packages/http-client/1.2.0/
\end{lstlisting}

The next time any project on your machine needs that same package at that same version, the package manager copies it from the cache instead of downloading it again.
You will see the message \texttt{(cached http-client@1.2.0)} in the install output when this happens.

\begin{tipbox}{Clearing the Cache}
If you ever need to force a fresh download, delete the cached package directory:

\begin{lstlisting}[language=bash, numbers=none]
rm -rf ~/.lattice/packages/http-client/1.2.0
clat install
\end{lstlisting}

Or clear the entire cache:

\begin{lstlisting}[language=bash, numbers=none]
rm -rf ~/.lattice/packages
\end{lstlisting}
\end{tipbox}

\subsection{The \texttt{lat\_modules/} Directory}\label{subsec:lat-modules-dir}
\index{lat\_modules!structure}

After installation, your project's dependency tree lives in \filename{lat\_modules/}:

\begin{lstlisting}[language=bash, numbers=none]
my-project/
  lattice.toml
  lattice.lock
  main.lat
  lat_modules/
    http-client/
      lattice.toml
      main.lat
    json-schema/
      lattice.toml
      src/
        main.lat
    cli-args/
      lattice.toml
      main.lat
\end{lstlisting}

Each package directory contains at minimum a \filename{lattice.toml} and one or more \filename{.lat} source files.
The module resolution system (described in \cref{sec:module-resolution}) knows how to find the entry point for each package.

To use an installed package in your code, import it by its package name:

\begin{lstlisting}[language=Lattice, caption={Importing installed packages}]
import "http-client" as http
import { validate } from "json-schema"
import "cli-args" as args

let response = http.get("https://api.weather.gov/points/44.9,-93.2")
let data = validate(response.body, schema)
\end{lstlisting}

The import system tries the built-in modules first, then falls back to \filename{lat\_modules/}---exactly as described in \cref{sec:module-resolution}.

%% ───────────────────────────────────────────────────────────
\section{Dependency Graphs and Cycle Detection}\label{sec:dep-graph}
\index{dependency graph}\index{circular dependencies}

When packages depend on other packages, those dependencies form a directed graph.
Lattice builds this graph after installation and checks it for cycles.

\subsection{How the Graph Is Built}\label{subsec:graph-construction}

The dependency graph is constructed by reading each installed package's \filename{lattice.toml}:

\begin{enumerate}
  \item Your project is the root node.
  \item Each of your direct dependencies becomes a node with an edge from the root.
  \item For each dependency, the package manager reads \emph{its} \filename{lattice.toml} and adds edges for its dependencies.
  \item This continues via breadth-first search\index{breadth-first search} until all transitive dependencies have been discovered.
\end{enumerate}

For example, if your project depends on \texttt{alpha}, and \texttt{alpha} depends on \texttt{beta} and \texttt{gamma}, the graph looks like:

\begin{center}
\begin{tikzpicture}[
  node distance=1.8cm,
  every node/.style={draw, rounded corners, fill=purple!5, minimum width=2cm, minimum height=0.7cm, font=\small},
  arrow/.style={-{Stealth[length=3mm]}, thick, color=latticepurple}
]
  \node (root) {my-project};
  \node (alpha) [below=of root] {alpha};
  \node (beta) [below left=of alpha] {beta};
  \node (gamma) [below right=of alpha] {gamma};

  \draw[arrow] (root) -- (alpha);
  \draw[arrow] (alpha) -- (beta);
  \draw[arrow] (alpha) -- (gamma);
\end{tikzpicture}
\end{center}

\subsection{Cycle Detection}\label{subsec:cycle-detection}
\index{circular dependencies!detection}

Circular dependencies---where \texttt{A} depends on \texttt{B}, \texttt{B} depends on \texttt{C}, and \texttt{C} depends on \texttt{A}---are a serious problem.
They make it impossible to determine a safe installation and loading order.

Lattice detects cycles using a depth-first search (DFS)\index{depth-first search} algorithm.
Each node is colored during traversal:

\begin{itemize}[nosep]
  \item \textbf{White}: not yet visited.
  \item \textbf{Gray}: currently being explored (on the DFS stack).
  \item \textbf{Black}: fully explored (all descendants processed).
\end{itemize}

If the DFS encounters a gray node---meaning we have followed a path back to a node that is still being explored---a cycle has been found.
The algorithm then traces back through the DFS stack to produce a human-readable cycle path:

\begin{lstlisting}[language=bash, numbers=none]
error: circular dependency detected: alpha -> beta -> gamma -> alpha
\end{lstlisting}

You can see this implementation in \filename{src/package.c}, in the \lstinline|dfs_visit()| and \lstinline|pkg_dep_graph_has_cycle()| functions.

\begin{warnbox}{Circular Dependencies Are Fatal}
If a circular dependency is detected, \texttt{clat install} fails and does not write a lock file.
You must resolve the cycle---typically by refactoring one of the packages to remove the circular relationship---before installation can succeed.
\end{warnbox}

%% ───────────────────────────────────────────────────────────
\section{Under the Hood: The Install Pipeline}\label{sec:install-pipeline}
\index{clat install!internals}

For readers who want the full picture, here is what happens when you run \texttt{clat install}, step by step.

\begin{enumerate}
  \item \textbf{Read \filename{lattice.toml}.}
  Parse the TOML into a \lstinline|PkgManifest| structure containing the project metadata and dependency list.
  If the file does not exist, print an error suggesting \texttt{clat init}.

  \item \textbf{Check for an existing lock file.}
  If \filename{lattice.lock} exists, parse it into a \lstinline|PkgLock| structure.
  The locked versions will be preferred during resolution.

  \item \textbf{Create \filename{lat\_modules/}.}
  If the directory does not exist, create it.

  \item \textbf{Resolve each dependency.}
  For each entry in \texttt{[dependencies]}:
  \begin{itemize}[nosep]
    \item If the lock file has an entry for this package, and the locked version satisfies the manifest constraint, use the locked version.
    \item Check if \filename{lat\_modules/\textlangle{}name\textrangle{}} already exists. If so, verify its version.
    \item If not installed: try a \lstinline|file://| registry (if configured), then fall back to the HTTP registry.
    \item On the HTTP registry: query available versions, select the best match, check the global cache (\filename{\textasciitilde/.lattice/packages/}), download if needed, copy to \filename{lat\_modules/}.
  \end{itemize}

  \item \textbf{Build the dependency graph.}
  Read each installed package's \filename{lattice.toml} to discover transitive dependencies.
  Construct a directed graph using BFS\@.

  \item \textbf{Check for cycles.}
  Run DFS-based cycle detection on the graph.
  If a cycle is found, report it and abort.

  \item \textbf{Write \filename{lattice.lock}.}
  Serialize the resolved versions, sources, and checksums to TOML\@.
\end{enumerate}

\begin{notebox}{Flat Installation}
\index{flat installation}
Lattice uses a \emph{flat} dependency layout: all packages live directly under \filename{lat\_modules/}, not nested inside each other.
This means that if two packages depend on the same third package, there is only one copy.
The trade-off is that version conflicts between transitive dependencies are not yet automatically resolved---this is a known limitation that will be addressed in a future version of the package manager.
\end{notebox}

%% ───────────────────────────────────────────────────────────
\section{Putting It All Together}\label{sec:package-workflow}

Let's walk through a complete workflow from project creation to running with dependencies.

\begin{lstlisting}[language=bash, numbers=none, caption={Complete project workflow}]
# Create a new project
mkdir sensor-dashboard && cd sensor-dashboard
clat init

# Add dependencies
clat add json-utils ^1.0.0
clat add http-server ^0.5.0
clat add template-engine ~2.1.0
\end{lstlisting}

Your \filename{lattice.toml} now looks like:

\begin{lstlisting}[language=bash, numbers=none]
[package]
name = "sensor-dashboard"
version = "0.1.0"

[dependencies]
json-utils = "^1.0.0"
http-server = "^0.5.0"
template-engine = "~2.1.0"
\end{lstlisting}

Write your application:

\begin{lstlisting}[language=Lattice, caption={main.lat --- using installed packages}]
import "json-utils" as json
import "http-server" as server
import { render } from "template-engine"

import "sensors/temperature" as temp
import "sensors/humidity" as humid

fn handle_dashboard(req: any) -> any {
    let readings = [
        temp.read_sensor(),
        humid.read_sensor()
    ]
    let body = render("dashboard.html", json.stringify(readings))
    return server.response(200, body)
}

let app = server.create(8080)
server.route(app, "/", handle_dashboard)
server.start(app)
print("Dashboard running on http://localhost:8080")
\end{lstlisting}

Notice how third-party packages (from \filename{lat\_modules/}) and your own modules (from \filename{sensors/}) are imported with the same syntax.
The module system handles the resolution transparently.

Run the program:

\begin{lstlisting}[language=bash, numbers=none]
clat main.lat
\end{lstlisting}

Share the project with a colleague:

\begin{lstlisting}[language=bash, numbers=none]
# In .gitignore:
lat_modules/

# Commit everything else:
git add lattice.toml lattice.lock main.lat sensors/
git commit -m "Initial sensor dashboard"
git push
\end{lstlisting}

When your colleague clones the repository, they run:

\begin{lstlisting}[language=bash, numbers=none]
clat install
clat main.lat
\end{lstlisting}

The lock file ensures they get the exact same dependency versions you used.

%% ───────────────────────────────────────────────────────────
\section{Exercises}\label{sec:ch29-exercises}

\begin{enumerate}
  \item \textbf{Initialize and explore.}
  Run \lstinline|clat init| in a new directory.
  Examine the generated \filename{lattice.toml}.
  Try running \lstinline|clat init| again---what happens?
  Now manually add a \texttt{description} and \texttt{license} field to the manifest.

  \item \textbf{Semver reasoning.}
  Given the available versions \texttt{[1.0.0, 1.1.0, 1.2.3, 1.3.0, 2.0.0, 2.1.0]}, which version does each constraint select?
  \begin{itemize}[nosep]
    \item \lstinline|"*"|
    \item \lstinline|"^1.1.0"|
    \item \lstinline|"~1.2.0"|
    \item \lstinline|">=1.3.0"|
    \item \lstinline|"<=1.2.3"|
    \item \lstinline|"1.2.3"|
  \end{itemize}

  \item \textbf{Local file registry.}
  Create a directory called \filename{my-registry} containing a subdirectory \filename{greetings} with a \filename{main.lat} that exports a \lstinline|hello| function.
  Set \lstinline|LATTICE_REGISTRY=file:///path/to/my-registry| and run \lstinline|clat add greetings|.
  Verify that the package appears in \filename{lat\_modules/} and can be imported.

  \item \textbf{Lock file investigation.}
  Add a dependency with a caret constraint.
  Examine the generated \filename{lattice.lock}.
  Change the version constraint in \filename{lattice.toml} to a tilde constraint for the same minor version.
  Run \lstinline|clat install| again.
  Does the lock file change?
  Why or why not?

  \item \textbf{Circular dependency detection.}
  Create three packages (\filename{alpha}, \filename{beta}, \filename{gamma}) in a local file registry.
  Make \texttt{alpha} depend on \texttt{beta}, \texttt{beta} depend on \texttt{gamma}, and \texttt{gamma} depend on \texttt{alpha}.
  Create a project that depends on \texttt{alpha} and run \lstinline|clat install|.
  Observe the error message.
  Then fix the cycle by removing one of the dependency edges and verify that the install succeeds.
\end{enumerate}

%% ───────────────────────────────────────────────────────────
\section*{What's Next}

With modules, imports, and packages under your belt, you have everything you need to build and share Lattice projects of any size.
But writing code is only half the story---maintaining it is the other half.
In \cref{ch:formatter-docgen}, we will explore Lattice's built-in code formatter and documentation generator, tools that help keep your codebase consistent and well-documented as it grows.
